\section{One Pell certificate for all positive powers}
\label{sec:batched-powers}

The principal new compression is at the level of existential certificates.
Every requested power retains its own input and output, but the auxiliary
number-theoretic environment is shared by the entire finite computation.
This is stronger than reusing identical arithmetic subexpressions: it
restricts independently chosen old witnesses to a common family, whose
simultaneous completeness must be proved.

For an integer $a\geq2$, let $\chi_a(t),\psi_a(t)\in\mathbb N$ be the Pell
coordinates defined by
\[
 \chi_a(t)+\psi_a(t)\sqrt{a^2-1}
       =(a+\sqrt{a^2-1})^t.
\]
These functions occur only in the proof. The emitted equation contains
ordinary polynomial expressions and no Pell-function symbol. Write
\[
 \mathcal N(a,x,y):\quad x^2+y^2=1+a^2y^2.
\]
All variables in this section range over $\mathbb N=\{0,1,2,\ldots\}$.

\subsection{The two established interfaces}

We use the positive-index case of the classical Pell-index representation.
For $a\geq2$ and $k>0$, the conjunction
\begin{align}
 &k\leq Y,\qquad \mathcal N(a,X,Y),\qquad
   \mathcal N(a,U,V),\qquad \mathcal N(B,S,R),\label{eq:index-norms}\\
 &B\geq2,\quad B\equiv1\pmod{4Y},\quad B\equiv a\pmod U,
   \quad V>0,\quad Y^2\mid V,\label{eq:index-frame}\\
 &S\equiv X\pmod U,\qquad R\equiv k\pmod{4Y}
 \label{eq:index-congruences}
\end{align}
is equivalent to $X=\chi_a(k)$ and $Y=\psi_a(k)$. In particular the
condition $k\leq Y$ makes $Y>0$ before any congruence modulus is used.

The positive-power representation says that, for $n,k>0$, $m=n^k$ is
equivalent to the existence of $W,a,T,Z,X,Y$ with $a\geq2$ satisfying
\begin{align}
 &X=\chi_a(k),\qquad Y=\psi_a(k),\label{eq:power-pell}\\
 &X+Yn\equiv Ya+m\pmod T,\qquad 2an=T+n^2+1,\label{eq:power-local}\\
 &m<T,\qquad n\leq W,\qquad k\leq W,
   \qquad \mathcal N(W+1,a,WZ).\label{eq:power-global}
\end{align}
These are the interfaces formalized as \texttt{Pell.matiyasevic} and
\texttt{Pell.eq\_pow\_of\_pell} in Mathlib's
\texttt{PellMatiyasevic.lean}~\cite{pell}. The shifted congruence in
\eqref{eq:power-local} avoids truncated subtraction. The standard proof
derives $a\geq n$, so it is equivalent to the source formulation
$X\equiv Y(a-n)+m\pmod T$.

\subsection{The batched polynomial conditions}

Consider a fixed nonempty list of power requests
\[
 (n_i,k_i,m_i),\qquad 1\leq i\leq s,
\]
whose bases and exponents are positive. Group identical exponent
expressions into $e_1,\ldots,e_g$, writing $k_i=e_{j(i)}$. Equality of
expressions is used here, not agreement at sampled numerical inputs.
Let
\[
 W=1+\sum_{i=1}^s(n_i+k_i+m_i).
 \label{eq:shared-W}
\]
Duplicate expressions may be included just once, provided the resulting
natural sum remains strictly larger than every base, exponent, and output.
Introduce a single frame $A,Z,U,V,B$, one quadruple
$(X_j,Y_j,S_j,R_j)$ for each exponent group, and one local modulus $T_i$
for each power request. Put
\[
 P=\prod_{j=1}^gY_j.
\]
The global conditions are
\begin{align}
 &A\geq2,\quad \mathcal N(W+1,A,WZ),\quad
   \mathcal N(A,U,V),\label{eq:batch-global-norms}\\
 &V=P^2(d+1),\qquad B=1+4Pq,
   \qquad B\equiv A\pmod U.\label{eq:batch-pooled}
\end{align}
Here $d,q$ are natural witnesses. For each exponent group retain
\begin{align}
 &e_j\leq Y_j,\quad \mathcal N(A,X_j,Y_j),\quad
     \mathcal N(B,S_j,R_j),\label{eq:batch-group-norms}\\
 &S_j\equiv X_j\pmod U,
       \qquad R_j\equiv e_j\pmod{4Y_j}.\label{eq:batch-group-congruences}
\end{align}
For each request retain only
\begin{equation}
 2An_i=T_i+n_i^2+1,\qquad
 X_{j(i)}+Y_{j(i)}n_i
      \equiv Y_{j(i)}A+m_i\pmod{T_i}.
 \label{eq:batch-local}
\end{equation}
The separate output-bound, $W$-bound, $B$-bound, local $V$-divisibility,
and local $B$-congruence certificates have disappeared.

\begin{theorem}[Simultaneous power certificate]
\label{thm:simultaneous-powers}
For every finite nonempty list of positive bases and exponents, the
conditions \eqref{eq:batch-global-norms}--\eqref{eq:batch-local} have
natural witnesses if and only if $m_i=n_i^{k_i}$ for every $i$.
If at least one base is known to be at least two, the explicit guard
$A\geq2$ may be omitted.
\end{theorem}

\begin{proof}[Soundness]
First suppose $A\geq2$. The last norm excludes $Z=0$, since that would
give $A=1$. Thus $Z\geq1$, and
\[
 A^2=1+(W^2+2W)W^2Z^2>W^2.
\]
Consequently $A>W$, which bounds every input and proposed output. Each
$Y_j\geq e_j>0$, so $P>0$ and
\[
 V=P^2(d+1)\geq Y_j^2\geq Y_j>0.
\]
In particular $Y_j^2\mid V$ and $B\equiv1\pmod{4Y_j}$.

The omitted guard $B\geq2$ follows from the retained equations. We have
$B=1+4Pq\geq1$. If $B=1$, then $\mathcal N(B,S_j,R_j)$ gives $S_j=1$.
The two norms with parameter $A$, together with $V\geq Y_j$, give
$U\geq X_j\geq2$. The congruence $S_j\equiv X_j\pmod U$ would then
imply $U\mid X_j-1$, although $0<X_j-1<U$. This is impossible.

All hypotheses of the Pell-index representation now hold group by group,
so $X_j=\chi_A(e_j)$ and $Y_j=\psi_A(e_j)$. For the omitted output bound,
$A\geq n_i+1$ yields
\[
 T_i-A=(2n_i-1)A-n_i^2-1
       \geq(n_i+2)(n_i-1)\geq0.
\]
Hence $m_i<W<A\leq T_i$. The standard positive-power representation
applies separately to every request and proves $m_i=n_i^{k_i}$.

Finally suppose the guard $A\geq2$ was omitted, but some $n_i\geq2$.
The retained modulus definition gives
$2An_i\geq n_i^2+1>2n_i$, and therefore $A>1$ before any preceding
argument is used. In the actual compiler a literal base-two request gives
the particularly simple equation $4A=T_i+5$.
\end{proof}

\begin{proof}[Simultaneous completeness]
Assume $m_i=n_i^{k_i}$ for every request, and first form the finite natural
number $W$ from these semantic inputs and outputs. Choose
\[
 A=\chi_{W+1}(W),\qquad
 Z=\frac{\psi_{W+1}(W)}{W}.
\]
The quotient is an integer because
$\psi_{W+1}(t)\equiv t\pmod W$. Also $W>0$, $A\geq2$, and
$\mathcal N(W+1,A,WZ)$ holds. For each $i$,
\[
 n_i^{k_i}\leq n_i^W<(W+1)^W\leq A.
\]
The standard Pell power congruence therefore supplies a positive
$T_i=2An_i-n_i^2-1>m_i$ and the local congruence
\eqref{eq:batch-local}. This is the source proof of the power
representation with a common upper bound in place of $\max(n_i,k_i)$.

Set $X_j=\chi_A(e_j)$ and $Y_j=\psi_A(e_j)$, so all $Y_j>0$, and form
$P=\prod_jY_j$. In the finite ring $\mathbb Z/(4P^2)$, the matrix
\[
 M=\begin{pmatrix}A&A^2-1\\1&A\end{pmatrix}
\]
is invertible because its determinant is one. Its powers belong to a
finite group. Some positive power is therefore the identity. Taking a
sufficiently large multiple $L$ of that period ensures $L>\max_j e_j$
and $L>1$. Put $U=\chi_A(L)$ and $V=\psi_A(L)$. The Pell matrix formula
gives
\[
 U\equiv1\pmod{4P^2},\qquad V\equiv0\pmod{4P^2}.
\]
Thus $V>0$, $P^2\mid V$, $U>A$, and $\gcd(U,4P)=1$.
The Chinese remainder theorem gives $B\geq2$ with
\[
 B\equiv A\pmod U,\qquad B\equiv1\pmod{4P}.
\]
Indeed the least nonnegative solution is already at least $A$, since
$A<U$. Define $d=V/P^2-1$ and $q=(B-1)/(4P)$.

For each group choose $S_j=\chi_B(e_j)$ and $R_j=\psi_B(e_j)$.
The recurrence with parameters congruent modulo $U$ yields
$S_j\equiv X_j\pmod U$. The recurrence modulo $B-1$ yields
$R_j\equiv e_j\pmod{4Y_j}$. All norms and bounds hold by construction,
so every retained condition is satisfied simultaneously.
\end{proof}

The finite-group step is deliberately stated for modulus $4P^2$.
Separately proving $Y_j^2\mid V$ for every $j$ would not prove
$P^2\mid V$, because the $Y_j$ can share prime factors. The common
period argument supplies the stronger divisibility actually used by the
emitted certificate.

\subsection{Compatibility with a common binomial radix}

The direct compiler also replaces the separate digit-extraction radices of
its finitely many binomial graphs by one radix. Given the positive upper arguments at the actual call sites,
$a_1,\ldots,a_r\in\mathbb N$, choose
\[
 H=\sum_{\ell=1}^{r}a_\ell,\qquad Q=2^H+1.
\]
Then $Q>2^{a_\ell}$ for every $\ell$. Every coefficient in
$(Q+1)^{a_\ell}$ is smaller than $Q$, since
$\binom{a_\ell}{j}\leq2^{a_\ell}$. Consequently quotient and remainder
equations for the digit at position $b_\ell$ recover
$\binom{a_\ell}{b_\ell}$ exactly, including the value zero when
$b_\ell>a_\ell$. Repeated lower arguments share the same power
$Q^{b_\ell}$. This is an ordinary positional-notation identity, with no
probabilistic or approximate digit extraction.

This common-radix construction and Theorem~\ref{thm:simultaneous-powers}
are compatible in a definite order. First choose a semantic solution of
the parent certificate and the common radix. Its finite power-request
list then has definite natural inputs and correct outputs. Next form
$W$ including all those outputs and choose the common Pell frame.
No Pell-frame coordinate enters the definition of the common radix or
the list of semantic power requests. Conversely, any zero of the final
equation yields all these natural values at once; the batched theorem
forces every requested power to have its intended value before the
binomial digit argument is invoked.

Both implications concern all natural numbers, including arbitrarily
large witnesses. A search over a finite numerical range is neither used
nor offered as a substitute for these proofs. The argument consists of
finite arithmetic, elementary Pell facts, finite-group periodicity, and
the Chinese remainder theorem, all formalizable in ZFC. Composing it
with the parent's direct C222 calibration preserves the stated
ZFC-provable consistency equivalence; it does not establish a new
kernel-checked or object-language ZFC derivation by itself.

\subsection{Eliminating terminal quadratic coordinates}

The final simplification acts on the literal quadratic residual system,
independently of the Pell implementation. Suppose one residual is
\[
 v-f(\mathbf x)=0,
\]
where $f\in\mathbb N[\mathbf x]$ has degree at most two and does not
contain $v$. If $f$ is affine, substitute it for $v$ everywhere. If $f$
is quadratic, make the same substitution provided $v$ occurs only
linearly in every other residual. Every new residual still has degree
at most two. Deleting the defining row and the variable $v$ gives an
exact bijection of natural zero sets: projection in one direction and
the unique restoration $v=f(\mathbf x)\in\mathbb N$ in the other.
The nonnegative coefficients are a sufficient, explicitly checked
condition for preserving the natural domain; a signed linear-algebra
elimination would not have this property automatically.

Repeated application provides a compact certificate consisting of the
deleted row, variable, and restoration polynomial at each step. Reading
the certificate backwards reconstructs every removed coordinate. The
verification checks exact polynomial identities, the domain condition,
and degree preservation at every step. It is stronger than comparing
the final polynomials at a selection of sample inputs.
